\documentclass[aspectratio=169,10pt]{beamer}

% Shared Beamer setup for Fall 2026 course slides.
% Clean Cal Poly-inspired theme: no custom class files or uncommon packages.
\mode<presentation>{
  \usetheme{default}
  \usefonttheme{professionalfonts}
}

\definecolor{CalPolyGreen}{HTML}{154734}
\definecolor{CalPolyGold}{HTML}{BD8B13}
\definecolor{CalPolySage}{HTML}{7FA28F}
\definecolor{CalPolyMint}{HTML}{DDF1E7}
\definecolor{CalPolySand}{HTML}{A29F86}
\definecolor{CalPolyBeige}{HTML}{EFEEDE}
\definecolor{DeckInk}{HTML}{1F2933}
\definecolor{DeckMuted}{HTML}{5F6B7A}
\definecolor{DeckSoft}{HTML}{F7F7F5}

\setbeamersize{text margin left=0.55in,text margin right=0.55in}
\setbeamercolor{normal text}{fg=DeckInk,bg=white}
\setbeamercolor{structure}{fg=CalPolyGreen}
\setbeamercolor{title}{fg=white}
\setbeamercolor{frametitle}{fg=CalPolyGreen,bg=white}
\setbeamercolor{block title}{bg=CalPolySage,fg=white}
\setbeamercolor{block body}{bg=CalPolyMint,fg=black}
\setbeamercolor{alerted text}{fg=CalPolyGold}

\setbeamerfont{title}{size=\Huge,series=\bfseries}
\setbeamerfont{frametitle}{size=\Large,series=\bfseries}
\setbeamerfont{block title}{size=\normalsize,series=\bfseries}

\setbeamertemplate{navigation symbols}{}
\setbeamertemplate{itemize item}{\color{CalPolyGold}$\blacktriangleright$}
\setbeamertemplate{itemize subitem}{\color{CalPolyGreen}$\bullet$}
\setbeamertemplate{footline}{
  \hfill{\color{DeckMuted}\scriptsize\insertframenumber/\inserttotalframenumber}
  \hspace{0.18in}\vspace{0.12in}
}
\setbeamertemplate{frametitle}{
  \vspace{0.18in}
  \hspace{0.02in}\insertframetitle\par
  \vspace{0.08in}
  \noindent\makebox[\linewidth][l]{%
    {\color{CalPolyGold}\rule{0.55in}{2pt}}%
    {\color{CalPolyGreen}\rule{\dimexpr\linewidth-0.55in\relax}{2pt}}%
  }
}

\newcommand{\CourseNumber}{Course}
\newcommand{\CourseTitle}{Course Title}
\newcommand{\LectureNumber}{0}
\newcommand{\LectureTitle}{Lecture Title}
\newcommand{\LectureDate}{Date}
\newcommand{\CourseUnit}{Unit}

\newcommand{\coursenumber}[1]{\renewcommand{\CourseNumber}{#1}}
\newcommand{\coursetitle}[1]{\renewcommand{\CourseTitle}{#1}}
\newcommand{\lecturenumber}[1]{\renewcommand{\LectureNumber}{#1}}
\newcommand{\lecturetitle}[1]{\renewcommand{\LectureTitle}{#1}}
\newcommand{\lecturedate}[1]{\renewcommand{\LectureDate}{#1}}
\newcommand{\courseunit}[1]{\renewcommand{\CourseUnit}{#1}}

\author{Kirk Duran}
\institute{California Polytechnic State University, San Luis Obispo}
\date{\LectureDate}

\newcommand{\makedecktitle}{
  \begingroup
  \setbeamercolor{background canvas}{bg=CalPolyGreen}
  \begin{frame}[plain]
    \vfill
    \begin{center}
      {\usebeamerfont{title}\color{white}\LectureTitle\par}
      \vspace{0.18in}
      {\Large\color{white}Lecture \LectureNumber\par}
    \end{center}
    \vfill
    \noindent{\color{white}\rule{\linewidth}{1.2pt}}\par
    \vspace{0.08in}
    \noindent\colorbox{CalPolyGold}{\makebox[0.24\linewidth][c]{\color{white}\Large\CourseNumber}}
    \hspace{0.12in}{\color{white}\Large Kirk Duran}
  \end{frame}
  \endgroup
}

\newenvironment{keyidea}{\begin{block}{Key idea}}{\end{block}}
\newenvironment{activity}{\begin{block}{In class}}{\end{block}}
\newenvironment{cpdefinition}{\begin{block}{Definition}}{\end{block}}
\newenvironment{exampleblockcp}{
  \begingroup
  \setbeamercolor{block title}{bg=CalPolySand,fg=white}
  \setbeamercolor{block body}{bg=CalPolyBeige,fg=black}
  \begin{block}{Example}
}{
  \end{block}
  \endgroup
}

\newcommand{\imageplaceholder}[2][0.3in]{%
  \noindent\makebox[\linewidth][c]{%
    \fcolorbox{CalPolySage}{CalPolyBeige}{%
      \parbox[c][#1][c]{0.86\linewidth}{%
        \centering
        {\color{DeckMuted}\scriptsize Image placeholder: }%
        {\color{DeckInk}\scriptsize\ttfamily\detokenize{#2}}%
      }%
    }%
  }%
}

\newcommand{\sectionframe}[1]{
  \begingroup
  \setbeamercolor{background canvas}{bg=CalPolyGreen}
  \begin{frame}[plain]
    \vfill
    {\color{white}\LARGE\bfseries #1\par}
    \vspace{0.18in}
    {\color{CalPolyGold}\rule{0.22\paperwidth}{3pt}}
    {\color{white}\rule{0.62\paperwidth}{3pt}}
    \vfill
  \end{frame}
  \endgroup
}

\newcommand{\placeholdernote}{
  \begin{block}{Placeholder}
    Replace these bullets with the final lecture material, examples, and in-class activities.
  \end{block}
}

\usepackage[T1]{fontenc}
\usepackage{lmodern}
\pdfstringdefDisableCommands{\def\translate#1{#1}}

\graphicspath{{../assets/lecture-27/}}

% Diagrams are reproducible with ../assets/lecture-27/generate-diagrams.py.
\newcommand{\lectureimage}[2][1.75in]{%
  \includegraphics[width=\linewidth,height=#1,keepaspectratio]{#2}%
}

\setbeamersize{text margin left=0.42in,text margin right=0.42in}

\coursenumber{CS 1001}
\coursetitle{Introduction to Computing}
\lecturenumber{27}
\lecturetitle{Functions over Structured Data}
\lecturedate{Fall 2026}
\courseunit{10. Structured Data}

\begin{document}

\makedecktitle

\begin{frame}{Today}
  \begin{itemize}
    \item Design functions whose inputs and outputs are lists, nested lists, strings, and dictionaries.
    \item Decompose structured-data problems into smaller transformations and summaries.
    \item Use selected Python tools---\texttt{sum}, \texttt{max}, \texttt{enumerate}, \texttt{zip}, \texttt{map}, \texttt{filter}, \texttt{any}, and \texttt{all}---when they express the intended computation directly.
    \item Apply \texttt{while} to repetition whose number of iterations is not known in advance.
    \item Preserve the distinction between producing a new structure and mutating an existing one.
  \end{itemize}

  \begin{keyidea}
    Today is a design lecture: each problem begins as a prompt, and we choose a representation and control structure that match the data and the required result.
  \end{keyidea}
\end{frame}

\begin{frame}{The Boundary for This Lecture}
  \begin{columns}[T,onlytextwidth]
    \begin{column}{0.47\textwidth}
      \begin{block}{Today}
        \begin{itemize}
          \item functions over structured values
          \item decomposition and helper functions
          \item selected sequence built-ins
          \item \texttt{map} and \texttt{filter}
          \item \texttt{enumerate} and \texttt{zip}
          \item \texttt{while}
          \item Boolean aggregation with \texttt{any} / \texttt{all}
        \end{itemize}
      \end{block}
    \end{column}
    \begin{column}{0.47\textwidth}
      \begin{alertblock}{Not today}
        \begin{itemize}
          \item searching algorithms
          \item sorting algorithms
          \item algorithmic complexity
          \item object-oriented design
          \item classes or methods we define ourselves
          \item advanced library APIs
        \end{itemize}
      \end{alertblock}
    \end{column}
  \end{columns}
\end{frame}

\sectionframe{Part 1: Prompt, Design, Implement, Refine}

\begin{frame}{How Today's Lecture Works}
  For each problem we will follow the same sequence.

  \begin{center}
    \textbf{Prompt} $\longrightarrow$ \textbf{Data Contract} $\longrightarrow$ \textbf{Plan} $\longrightarrow$ \textbf{Code} $\longrightarrow$ \textbf{Refine}
  \end{center}

  \begin{itemize}
    \item First decide what the input and output structures mean.
    \item Then identify the repeated computation or changing state.
    \item Write a correct version using mechanisms we already understand.
    \item Finally ask whether a Python construct communicates that same intent more directly.
  \end{itemize}

  \begin{keyidea}
    A built-in is useful when it removes incidental mechanism without hiding the computation we need to understand.
  \end{keyidea}
\end{frame}

\begin{frame}{Structured Data Appears in the Function Contract}
  \small
  \begin{columns}[T,onlytextwidth]
    \begin{column}{0.55\textwidth}
      \begin{block}{Examples}
        \footnotesize
        \texttt{def average(values: list[float]) -> float:}\\[0.08in]
        \texttt{def flatten(grid: list[list[int]]) -> list[int]:}\\[0.08in]
        \texttt{def row\_maxima(grid: list[list[int]]) -> dict[int, int]:}
      \end{block}

      \begin{itemize}
        \item The annotation communicates more than element type.
        \item It communicates the \textbf{shape of the expected data}; annotations do not enforce types or row lengths at runtime.
        \item The return type states what representation the function promises to produce.
      \end{itemize}
    \end{column}
    \begin{column}{0.39\textwidth}
      % Intended visual: three function contracts connected to 1D list, 2D grid, and dictionary output.
      \lectureimage[2.05in]{structured-function-contracts.png}
    \end{column}
  \end{columns}
\end{frame}

\begin{frame}{Prompt 1: Compute an Average}
  \begin{activity}
    Write a function that accepts a nonempty list of numerical values and returns their arithmetic mean.
  \end{activity}

  \begin{block}{Contract}
    \texttt{def average(values: list[float]) -> float:}
  \end{block}

  Before writing code:
  \begin{enumerate}
    \item What quantity must be accumulated?
    \item What second quantity is needed to produce the mean?
    \item Does the input list need to be modified?
  \end{enumerate}
\end{frame}

\begin{frame}{A Correct Loop-Based Solution}
  \begin{block}{Mechanism made explicit}
    \texttt{def average(values: list[float]) -> float:}\\
    \hspace*{0.20in}\texttt{total = 0.0}\\
    \hspace*{0.20in}\texttt{for value in values:}\\
    \hspace*{0.40in}\texttt{total = total + value}\\
    \hspace*{0.20in}\texttt{return total / len(values)}
  \end{block}

  \begin{itemize}
    \item \texttt{total} is the result-so-far accumulator.
    \item \texttt{len(values)} supplies the number of observations.
    \item The input is read but never mutated.
  \end{itemize}
\end{frame}

\begin{frame}{When the Built-In Matches the Intent}
  Python already provides \texttt{sum} for numerical accumulation.

  \begin{block}{Refined implementation}
    \texttt{def average(values: list[float]) -> float:}\\
    \hspace*{0.20in}\texttt{return sum(values) / len(values)}
  \end{block}

  \begin{center}
    \textbf{explicit accumulation} $\longrightarrow$ \textbf{named aggregate operation}
  \end{center}

  \begin{keyidea}
    The function contract and mathematical meaning have not changed. Only the level at which the accumulation is expressed has changed.
  \end{keyidea}
\end{frame}

\begin{frame}{Prompt 2: Average of Averages}
  \begin{activity}
    A nested list stores several groups of scores. Return the average of the average score from each group.
  \end{activity}

  \begin{block}{Example input}
    \texttt{groups = [[80, 90], [70, 100], [85, 95]]}
  \end{block}

  Desired intermediate structure:
  \begin{center}
    \texttt{[85.0, 85.0, 90.0]}
  \end{center}

  Then average those three values. Assume the outer list and every group are nonempty.
\end{frame}

\begin{frame}{Decompose the Problem with a Helper}
  \begin{block}{Reuse the abstraction we just designed}
    \texttt{def average(values: list[float]) -> float:}\\
    \hspace*{0.20in}\texttt{return sum(values) / len(values)}\\[0.10in]
    \texttt{def average\_of\_averages(groups: list[list[float]]) -> float:}\\
    \hspace*{0.20in}\texttt{group\_averages = []}\\
    \hspace*{0.20in}\texttt{for group in groups:}\\
    \hspace*{0.40in}\texttt{group\_averages.append(average(group))}\\
    \hspace*{0.20in}\texttt{return average(group\_averages)}
  \end{block}

  \begin{keyidea}
    The outer function reasons about groups. The helper function owns the lower-level meaning of ``average one group.''
  \end{keyidea}
\end{frame}

\begin{frame}{Average of Averages Is a Specific Contract}
  Suppose groups have different lengths:

  \begin{block}{Data}
    \texttt{[[100], [0, 0, 0]]}
  \end{block}

  \begin{columns}[T,onlytextwidth]
    \begin{column}{0.47\textwidth}
      \begin{block}{Average of group averages}
        \centering
        \texttt{(100 + 0) / 2 = 50}
      \end{block}
      Each group receives equal weight.
    \end{column}
    \begin{column}{0.47\textwidth}
      \begin{block}{Average of all observations}
        \centering
        \texttt{(100 + 0 + 0 + 0) / 4 = 25}
      \end{block}
      Each observation receives equal weight.
    \end{column}
  \end{columns}

  \begin{alertblock}{Design consequence}
    A function name and purpose statement must say which quantity is intended.
  \end{alertblock}
\end{frame}

\sectionframe{Part 2: Functions over Nested Structures}

\begin{frame}{Prompt 3: Flatten One Level}
  \begin{activity}
    Write a function that accepts a two-dimensional list and returns a new one-dimensional list containing the elements row by row.
  \end{activity}

  \begin{block}{Example}
    \texttt{[[1, 2], [3, 4, 5], [6]]}\\
    $\longrightarrow$\\
    \texttt{[1, 2, 3, 4, 5, 6]}
  \end{block}

  Questions:
  \begin{enumerate}
    \item How many structural levels must be traversed?
    \item What state should accumulate the result?
    \item Should any row in the input be modified?
  \end{enumerate}
\end{frame}

\begin{frame}{Flattening Requires Traversing Two Structural Levels}
  \begin{columns}[T,onlytextwidth]
    \begin{column}{0.54\textwidth}
      \begin{center}
        \textbf{outer list} $\rightarrow$ \textbf{row} $\rightarrow$ \textbf{element}
      \end{center}

      \begin{itemize}
        \item The outer loop selects one row at a time.
        \item The inner loop selects one element from that row at a time.
        \item A fresh output list accumulates the elements in encounter order.
      \end{itemize}

      \begin{keyidea}
        Nested data usually leads to nested processing because each loop traverses one structural level.
      \end{keyidea}
    \end{column}
    \begin{column}{0.40\textwidth}
      % Intended visual: rows of a grid flowing left-to-right into one flat sequence.
      \lectureimage[2.50in]{flatten-grid-to-sequence.png}
    \end{column}
  \end{columns}
\end{frame}

\begin{frame}{Implementation: Flatten One Level}
  \begin{block}{Function}
    \texttt{def flatten(grid: list[list[int]]) -> list[int]:}\\
    \hspace*{0.20in}\texttt{result = []}\\
    \hspace*{0.20in}\texttt{for row in grid:}\\
    \hspace*{0.40in}\texttt{for value in row:}\\
    \hspace*{0.60in}\texttt{result.append(value)}\\
    \hspace*{0.20in}\texttt{return result}
  \end{block}

  \begin{itemize}
    \item \texttt{grid} and its rows are read only.
    \item \texttt{result} is intentionally mutated because it is a new list owned by this function.
    \item The returned structure is independent of the outer input list.
  \end{itemize}
\end{frame}

\begin{frame}{Prompt 4: Record the Maximum of Each Row}
  \begin{activity}
    Given a two-dimensional list of integers, return a dictionary mapping each row index to the maximum value in that row.
  \end{activity}

  \begin{block}{Example}
    \texttt{[[4, 9, 2], [7, 1, 8], [5, 3, 6]]}\\
    $\longrightarrow$\\
    \texttt{\{0: 9, 1: 8, 2: 6\}}
  \end{block}

  Assume every row is nonempty, as required by \texttt{max}. This problem needs both:
  \begin{itemize}
    \item each row's \textbf{position}, and
    \item the row \textbf{value} itself.
  \end{itemize}
\end{frame}

\begin{frame}{\texttt{enumerate}: Position and Value Together}
  Compare:

  \begin{columns}[T,onlytextwidth]
    \begin{column}{0.47\textwidth}
      \begin{block}{Index-driven}
        \texttt{for i in range(len(rows)):}\\
        \hspace*{0.20in}\texttt{row = rows[i]}
      \end{block}
    \end{column}
    \begin{column}{0.47\textwidth}
      \begin{block}{Enumerated}
        \texttt{for i, row in enumerate(rows):}\\
        \hspace*{0.20in}\texttt{...}
      \end{block}
    \end{column}
  \end{columns}

  \begin{cpdefinition}
    \texttt{enumerate(sequence)} supplies pairs containing an element's index and the element value.
  \end{cpdefinition}

  Use it when the computation genuinely needs both pieces of information.
\end{frame}

\begin{frame}{Implementation: Row Maxima}
  \begin{block}{Function}
    \texttt{def row\_maxima(grid: list[list[int]]) -> dict[int, int]:}\\
    \hspace*{0.20in}\texttt{result = \{\}}\\
    \hspace*{0.20in}\texttt{for index, row in enumerate(grid):}\\
    \hspace*{0.40in}\texttt{result[index] = max(row)}\\
    \hspace*{0.20in}\texttt{return result}
  \end{block}

  \begin{center}
    \textbf{row position} $+$ \textbf{row aggregate} $\longrightarrow$ \textbf{dictionary association}
  \end{center}

  \begin{keyidea}
    The output representation is part of the solution: the dictionary preserves an explicit association between each row index and its summary value.
  \end{keyidea}
\end{frame}

\sectionframe{Part 3: Named Transformation Patterns --- \texttt{map} and \texttt{filter}}

\begin{frame}{Prompt 5: Apply One Function to Every Value}
  \begin{activity}
    Given a list of integers, produce a new list containing the square of each value.
  \end{activity}

  We already know two representations:

  \begin{block}{Explicit loop}
    \texttt{result = []}\\
    \texttt{for value in values:}\\
    \hspace*{0.20in}\texttt{result.append(square(value))}
  \end{block}

  \begin{block}{Comprehension}
    \texttt{result = [square(value) for value in values]}
  \end{block}
\end{frame}

\begin{frame}{\texttt{map}: Name the Transformation Pattern Directly}
  \begin{block}{Helper function}
    \texttt{def square(value: int) -> int:}\\
    \hspace*{0.20in}\texttt{return value * value}
  \end{block}

  \begin{block}{Mapping}
    \texttt{result = list(map(square, values))}
  \end{block}

  \begin{cpdefinition}
    \texttt{map(function, sequence)} applies the function to each supplied value. In Python 3, \texttt{map} produces an iterator: values are generated as it is consumed. \texttt{list(...)} collects them into a list.
  \end{cpdefinition}

  \begin{center}
    \textbf{same computational pattern} --- different representation
  \end{center}
\end{frame}

\begin{frame}{\texttt{filter}: Name the Selection Pattern Directly}
  \begin{block}{Predicate}
    \texttt{def positive(value: int) -> bool:}\\
    \hspace*{0.20in}\texttt{return value > 0}
  \end{block}

  \begin{block}{Filtering}
    \texttt{result = list(filter(positive, values))}
  \end{block}

  \begin{cpdefinition}
    \texttt{filter(predicate, sequence)} retains the values for which the predicate returns \texttt{True}.
  \end{cpdefinition}

  The predicate decides \textbf{whether} an existing value remains; it does not transform that value. Like \texttt{map}, \texttt{filter} returns an iterator.
\end{frame}

\begin{frame}{Loop, Comprehension, \texttt{map}, or \texttt{filter}?}
  \begin{center}
    \renewcommand{\arraystretch}{1.25}
    \begin{tabular}{l|l}
      \textbf{Representation} & \textbf{Good fit} \\\hline
      explicit \texttt{for} & multiple statements or changing state \\
      comprehension & compact transform / select into a new list \\
      \texttt{map} & apply an existing named function to every value \\
      \texttt{filter} & retain values using an existing named predicate
    \end{tabular}
  \end{center}

  \begin{keyidea}
    These forms are not ranked from ``basic'' to ``advanced.'' Choose the representation that makes the data relationship easiest to read.
  \end{keyidea}
\end{frame}

\begin{frame}{Prompt 6: Caesar Cipher}
  \begin{activity}
    Write a function that shifts each lowercase English letter forward by a fixed number of positions, wrapping from \texttt{z} back to \texttt{a}.
  \end{activity}

  \begin{block}{Example with shift 3}
    \texttt{"attack"} $\longrightarrow$ \texttt{"dwwdfn"}
  \end{block}

  Assume the input contains only lowercase letters \texttt{a} through \texttt{z}. This combines ideas from several units:
  \begin{itemize}
    \item strings as immutable sequences,
    \item \texttt{ord} / \texttt{chr},
    \item modular arithmetic,
    \item iteration and function decomposition.
  \end{itemize}
\end{frame}

\begin{frame}{Design a One-Character Transformation First}
  For lowercase letters, \texttt{ord('a')} provides a useful origin.

  \begin{block}{Helper}
    \texttt{def shift\_char(char: str, shift: int) -> str:}\\
    \hspace*{0.20in}\texttt{offset = ord(char) - ord("a")}\\
    \hspace*{0.20in}\texttt{shifted = (offset + shift) \% 26}\\
    \hspace*{0.20in}\texttt{return chr(ord("a") + shifted)}
  \end{block}

  \begin{center}
    character $\rightarrow$ numeric offset $\rightarrow$ shifted offset $\rightarrow$ character
  \end{center}
\end{frame}

\begin{frame}{Build the String from Transformed Characters}
  \begin{block}{Function}
    \texttt{def caesar(text: str, shift: int) -> str:}\\
    \hspace*{0.20in}\texttt{result = ""}\\
    \hspace*{0.20in}\texttt{for char in text:}\\
    \hspace*{0.40in}\texttt{result = result + shift\_char(char, shift)}\\
    \hspace*{0.20in}\texttt{return result}
  \end{block}

  \begin{keyidea}
    Because strings are immutable, each update rebinds \texttt{result} to a newly produced string. The input \texttt{text} is never mutated.
  \end{keyidea}
\end{frame}

\sectionframe{Part 4: Functions over Corresponding Structures}

\begin{frame}{Prompt 7: Match Student Names with Grades}
  \begin{activity}
    Two lists store corresponding data: one contains names and the other contains scores. Return the names of students whose score is at least 70.
  \end{activity}

  \begin{block}{Data}
    \texttt{names  = ["Ada", "Grace", "Alan", "Edsger"]}\\
    \texttt{scores = [92, 68, 81, 65]}
  \end{block}

  Desired result:
  \begin{center}
    \texttt{["Ada", "Alan"]}
  \end{center}
\end{frame}

\begin{frame}{\texttt{zip}: Corresponding Values Together}
  \begin{columns}[T,onlytextwidth]
    \begin{column}{0.54\textwidth}
      \begin{block}{Iteration}
        \texttt{for name, score in zip(names, scores):}\\
        \hspace*{0.20in}\texttt{...}
      \end{block}

      \begin{cpdefinition}
        \texttt{zip} groups values from corresponding positions of multiple sequences.
      \end{cpdefinition}

      \begin{itemize}
        \item first name with first score,
        \item second name with second score,
        \item and so on.
      \end{itemize}
    \end{column}
    \begin{column}{0.40\textwidth}
      % Intended visual: two parallel sequences connected pairwise into tuples.
      \lectureimage[2.4in]{zip-corresponding-sequences.png}
    \end{column}
  \end{columns}
\end{frame}

\begin{frame}{Implementation: Passing Students}
  \begin{block}{Function}
    \texttt{def passing\_students(}\\
    \hspace*{0.20in}\texttt{names: list[str], scores: list[int]}\\
    \texttt{) -> list[str]:}\\
    \hspace*{0.20in}\texttt{result = []}\\
    \hspace*{0.20in}\texttt{for name, score in zip(names, scores):}\\
    \hspace*{0.40in}\texttt{if score >= 70:}\\
    \hspace*{0.60in}\texttt{result.append(name)}\\
    \hspace*{0.20in}\texttt{return result}
  \end{block}

  \begin{alertblock}{Function contract}
    This design assumes equal lengths and corresponding data in the same order. By default, \texttt{zip} stops at the shorter input; check lengths when every entry must have a partner.
  \end{alertblock}
\end{frame}

\sectionframe{Part 5: When Repetition Depends on a Condition}

\begin{frame}{Prompt 8: Keep Asking Until the Answer Is Correct}
  \begin{activity}
    Ask the user \texttt{"What is 4 + 5?"}. If the answer is wrong, ask again. Stop only after the user enters \texttt{9}.
  \end{activity}

  Assume each response can be converted to an integer. The number of repetitions is unknown:

  \begin{itemize}
    \item the user may answer correctly immediately;
    \item the user may require several attempts;
    \item there is no input sequence whose length determines the repetition count.
  \end{itemize}

  \begin{keyidea}
    This is condition-controlled repetition rather than sequence-controlled iteration.
  \end{keyidea}
\end{frame}

\begin{frame}{The \texttt{while} Statement}
  \begin{columns}[T,onlytextwidth]
    \begin{column}{0.54\textwidth}
      \begin{block}{General form}
        \texttt{while condition:}\\
        \hspace*{0.20in}\texttt{body}
      \end{block}

      Execution model:
      \begin{enumerate}
        \item evaluate the condition;
        \item if \texttt{False}, leave the loop;
        \item if \texttt{True}, execute the body;
        \item return to the condition and evaluate it again.
      \end{enumerate}
    \end{column}
    \begin{column}{0.40\textwidth}
      % Intended visual: decision diamond feeding loop body and returning to the condition.
      \lectureimage[2.45in]{while-loop-control-flow.png}
    \end{column}
  \end{columns}
\end{frame}

\begin{frame}{Implementation: Repeated Math Prompt}
  \begin{block}{Function}
    \texttt{def ask\_addition() -> None:}\\
    \hspace*{0.20in}\texttt{answer = int(input("What is 4 + 5? "))}\\
    \hspace*{0.20in}\texttt{while answer != 9:}\\
    \hspace*{0.40in}\texttt{print("Try again.")}\\
    \hspace*{0.40in}\texttt{answer = int(input("What is 4 + 5? "))}\\
    \hspace*{0.20in}\texttt{print("Correct.")}
  \end{block}

  \begin{keyidea}
    Something in the loop body must make progress toward changing the condition. Here, a new user response replaces the previous value of \texttt{answer}.
  \end{keyidea}
\end{frame}

\begin{frame}{Choose \texttt{for} or \texttt{while} from the Source of Repetition}
  \begin{center}
    \renewcommand{\arraystretch}{1.30}
    \begin{tabular}{l|l|l}
      \textbf{Question} & \textbf{Typical construct} & \textbf{Example} \\\hline
      process each value in a sequence & \texttt{for} & each row in a grid \\
      repeat a known count & \texttt{for} + \texttt{range} & five trials \\
      repeat while a condition holds & \texttt{while} & ask until correct
    \end{tabular}
  \end{center}

  \begin{alertblock}{Common failure mode}
    A \texttt{while} loop can run forever if its condition never becomes \texttt{False}.
  \end{alertblock}
\end{frame}

\sectionframe{Part 6: Climax --- Password Validation with \texttt{any} and \texttt{all}}

\begin{frame}{Final Prompt: Validate a Password}
  \begin{activity}
    Write a function that returns \texttt{True} exactly when a password:
    \begin{itemize}
      \item has at least 8 characters,
      \item contains at least one uppercase letter,
      \item contains at least one lowercase letter,
      \item contains at least one digit, and
      \item contains no whitespace characters.
    \end{itemize}
  \end{activity}

  \begin{block}{Contract}
    \texttt{def valid\_password(password: str) -> bool:}
  \end{block}
\end{frame}

\begin{frame}{\texttt{any}: Does At Least One Result Satisfy the Condition?}
  Suppose we compute one Boolean per character:

  \begin{block}{Uppercase checks}
    \texttt{[char.isupper() for char in password]}
  \end{block}

  \begin{cpdefinition}
    \texttt{any(values)} returns \texttt{True} when at least one supplied value is truthy.
  \end{cpdefinition}

  Therefore:

  \begin{block}{Requirement}
    \texttt{has\_upper = any([char.isupper() for char in password])}
  \end{block}

  The same pattern works for lowercase letters and digits. \texttt{any([])} is \texttt{False}; here the comprehension builds its Boolean list before \texttt{any} combines the results.
\end{frame}

\begin{frame}{\texttt{all}: Does Every Result Satisfy the Condition?}
  For whitespace, the requirement applies to \textbf{every} character.

  \begin{block}{Per-character claim}
    \texttt{not char.isspace()}
  \end{block}

  \begin{cpdefinition}
    \texttt{all(values)} returns \texttt{True} when every supplied value is truthy.
  \end{cpdefinition}

  Therefore:

  \begin{block}{Requirement}
    \texttt{no\_space = all([not char.isspace() for char in password])}
  \end{block}

  \begin{center}
    \texttt{any}: existential condition \hspace{0.35in} \texttt{all}: universal condition\\[0.06in]
    \texttt{all([])} is \texttt{True}: no element violates the condition.
  \end{center}
\end{frame}

\begin{frame}{Build Each Requirement Independently}
  \begin{block}{Intermediate Boolean values}
    \texttt{long\_enough = len(password) >= 8}\\
    \texttt{has\_upper = any([char.isupper() for char in password])}\\
    \texttt{has\_lower = any([char.islower() for char in password])}\\
    \texttt{has\_digit = any([char.isdigit() for char in password])}\\
    \texttt{no\_space = all([not char.isspace() for char in password])}
  \end{block}

  \begin{keyidea}
    Instead of one complicated expression, name each independent requirement so its meaning can be inspected and tested separately.
  \end{keyidea}
\end{frame}

\begin{frame}[fragile]{Complete Password Checker}
  \small
  \begin{columns}[T,onlytextwidth]
    \begin{column}{0.60\textwidth}
{\footnotesize
\begin{verbatim}
def valid_password(password: str) -> bool:
    long_enough = len(password) >= 8
    has_upper = any([char.isupper()
                    for char in password])
    has_lower = any([char.islower()
                    for char in password])
    has_digit = any([char.isdigit()
                    for char in password])
    no_space = all([not char.isspace()
                    for char in password])
    return (long_enough and has_upper and
            has_lower and has_digit and no_space)
\end{verbatim}
}
    \end{column}
    \begin{column}{0.34\textwidth}
      \lectureimage[2.35in]{password-requirement-pipeline.png}
    \end{column}
  \end{columns}
\end{frame}

\begin{frame}{Trace the Requirements, Not Just the Final Boolean}
  \begin{center}
    \renewcommand{\arraystretch}{1.23}
    \begin{tabular}{l|c|c|c|c|c|c}
      \textbf{Password} & \textbf{Len} & \textbf{Upper} & \textbf{Lower} & \textbf{Digit} & \textbf{No space} & \textbf{Valid} \\\hline
      \texttt{CompSci9} & T & T & T & T & T & T \\
      \texttt{compsci9} & T & F & T & T & T & F \\
      \texttt{Comp Sci9} & T & T & T & T & F & F \\
      \texttt{CS9abc} & F & T & T & T & T & F
    \end{tabular}
  \end{center}

  \begin{keyidea}
    A complex validation function becomes easier to reason about when each rule is represented by one named Boolean claim.
  \end{keyidea}
\end{frame}

\begin{frame}{The Structured-Data Toolkit We Used Today}
  \begin{center}
    \renewcommand{\arraystretch}{1.25}
    \begin{tabular}{l|l}
      \textbf{Tool} & \textbf{Question it answers} \\\hline
      \texttt{sum} / \texttt{max} & how can this structure be summarized? \\
      \texttt{enumerate} & what are the position and value together? \\
      \texttt{zip} & which values from different structures correspond? \\
      \texttt{map} & what should every value become? \\
      \texttt{filter} & which existing values should remain? \\
      \texttt{any} & does at least one Boolean condition hold? \\
      \texttt{all} & do all Boolean conditions hold? \\
      \texttt{while} & should repetition continue based on current state?
    \end{tabular}
  \end{center}
\end{frame}

\begin{frame}{Takeaways}
  \begin{enumerate}
    \item The shape and meaning of structured data belong in a function's contract.
    \item Decompose a large structured-data problem into smaller transformations and summaries.
    \item Use nested iteration when the data itself has nested structure.
    \item Use \texttt{enumerate} when a computation needs position and value; use \texttt{zip} when it needs corresponding values from multiple sequences.
    \item \texttt{map} and \texttt{filter} are named forms of transformation and selection patterns already studied with loops and comprehensions.
    \item Use \texttt{while} when repetition is governed by a condition rather than by a sequence.
    \item \texttt{any} and \texttt{all} let many Boolean results collapse into one higher-level decision.
  \end{enumerate}

  \begin{block}{Next}
    The next unit introduces \textbf{objects and compound data}: a new way to group related state under a single programmer-defined abstraction.
  \end{block}
\end{frame}

\end{document}
